% ============================================================
% 10_math.tex — Math-Macros (generisch, fachunabhängig)
% ============================================================

% --- Zahlmengen (Blackboard Bold) ---
\newcommand{\R}{\mathbb{R}}
\newcommand{\C}{\mathbb{C}}
\newcommand{\N}{\mathbb{N}}
\newcommand{\Z}{\mathbb{Z}}
\newcommand{\Q}{\mathbb{Q}}

% --- Differentialoperator (aufrechtes d) ---
\newcommand{\dd}{\,\mathrm{d}}

% --- Operatornamen ---
\DeclareMathOperator{\sgn}{sgn}
\DeclareMathOperator{\grad}{grad}
\DeclareMathOperator{\divg}{div}
\DeclareMathOperator{\rot}{rot}
\newcommand{\vect}[1]{\vec{#1}}

% --- Gepaarte Begrenzer (mathtools) ---
\DeclarePairedDelimiter{\abs}{\lvert}{\rvert}
\DeclarePairedDelimiter{\norm}{\lVert}{\rVert}

% --- Summen-Mode-System: zentrales Umschalten side/stack/auto ---
\newcommand{\ZSFsumMode}{auto}
\newcommand{\SetZSFsumMode}[1]{\renewcommand{\ZSFsumMode}{#1}}
\newcommand{\ZSFsumSide}[2]{\sum\nolimits_{#1}^{#2}}
\newcommand{\ZSFsumStack}[2]{\sum\limits_{#1}^{#2}}
\newcommand{\ZSFsumAuto}[2]{%
  \IfStrEqCase{\ZSFsumMode}{%
    {side}{\ZSFsumSide{#1}{#2}}%
    {stack}{\ZSFsumStack{#1}{#2}}%
    {auto}{\sum_{#1}^{#2}}%
  }[\sum_{#1}^{#2}]%
}

% --- Limes-Mode-System (analog) ---
\newcommand{\ZSFlimMode}{auto}
\newcommand{\SetZSFlimMode}[1]{\renewcommand{\ZSFlimMode}{#1}}
\newcommand{\ZSFlimAuto}[1]{%
  \IfStrEqCase{\ZSFlimMode}{%
    {side}{\lim\nolimits_{#1}}%
    {stack}{\lim\limits_{#1}}%
    {auto}{\lim_{#1}}%
  }[\lim_{#1}]%
}

% --- Semantisches Formel-Highlighting (POSITIONAL) ---
% Verknüpft visuell zusammengehörige Teile einer Formel über mehrere
% Umformungsschritte hinweg. Farben zentral in 40_colors_structure.tex.
% Die Rolle ist die STELLE in der Herleitung, nicht die Grösse — für
% dokumentweite Grössen-Identität siehe \ZSFDeclareQuantity unten.
% Die Farbe läuft über \ZSFInk, damit sie auf einer Fläche, die ihre eigene
% Kontrastfarbe setzt, nicht unlesbar wird (40_colors_structure → Ink-Vertrag).
\newcommand{\ZSFmhlA}[1]{\ZSFInk{MathHLPrimary}{#1}}
\newcommand{\ZSFmhlB}[1]{\ZSFInk{MathHLSecondary}{#1}}
\newcommand{\ZSFmhlC}[1]{\ZSFInk{MathHLTarget}{#1}}
\newcommand{\ZSFmhlD}[1]{\ZSFInk{MathHLTertiary}{#1}}

% Für die dokumentweite Grössen-Identität (Farbe = Grösse statt Farbe =
% Stelle in der Herleitung) siehe \ZSFDeclareQuantity in
% styles/40_colors_structure.tex — dort, weil die Palette dort steht.

% --- Term-Annotation mit geschweifter Klammer (Stilmittel) ---
% Beschriftet einen Formelteil unter/über einer geschweiften Klammer.
% Das Label wird via \text gesetzt (aufrecht, kein Mathe-Spacing); für
% Mathe im Label $...$ verwenden, z.B. \ZSFbraceunder{x}{$=0$ im Vakuum}.
%   \[ \varphi = \frac{q}{r} - \ZSFbraceover{\frac{q}{r_B}}{Bezugspotenzial} \]
\newcommand{\ZSFbraceunder}[2]{\underbrace{#1}_{\text{#2}}}
\newcommand{\ZSFbraceover}[2]{\overbrace{#1}^{\text{#2}}}

% --- Lange Pfeile: bewusst auf die kurzen Varianten umgebogen ---
% Historisch ein sansmath-Bugfix (der \neq-Fix von damals ist mit unicode-math
% ersatzlos entfallen — die Schrift bringt das Zeichen selbst mit). Die
% Umbiegung der langen Pfeile bleibt als Platz-Entscheidung: In ~50 mm
% schmalen Spalten ist der kurze Pfeil die richtige Breite, und die Kapitel
% verlassen sich auf diese Masse.
\renewcommand{\implies}{\;\Rightarrow\;}
\renewcommand{\iff}{\;\Leftrightarrow\;}
\renewcommand{\Longrightarrow}{\Rightarrow}
\renewcommand{\Longleftrightarrow}{\Leftrightarrow}
\renewcommand{\longrightarrow}{\rightarrow}
\renewcommand{\longleftarrow}{\leftarrow}
